% ============================================================
\chapter{Pre-trained Models}
\label{ch:pretrained}
% ============================================================

In 2018, three papers transformed NLP: ELMo (Peters et al.), GPT (Radford et al.), and above all BERT (Devlin et al.). Their shared idea: pre-train a large language model on billions of unlabelled words, then fine-tune it on the target task with very little supervised data. This ``pre-training / fine-tuning'' paradigm revolutionised the field: the learned representations capture syntax, semantics, and even factual knowledge about the world, rendering previous specialised approaches largely obsolete. This was the beginning of the \emph{foundation model} era.

\begin{intuition}[title=\textbf{The pre-training / fine-tuning paradigm}]
Rather than training a model from scratch for each task, we pre-train a large
model on unlabeled data, then adapt (fine-tune) it on task-specific data. This
paradigm, inaugurated by ELMo (2018) and popularized by BERT (2018), has
transformed NLP.
\end{intuition}

% ------------------------------------------------------------
\section{Learning Contextual Representations}
\label{sec:pre:contextual}
% ------------------------------------------------------------

\begin{definition}[Contextual Representation]
A \emph{contextual representation} is a function
$\phi : \mathcal{V}^T \times \{1, \ldots, T\} \to \R^d$ that maps each token
$x_t$ to a vector $\mathbf{h}_t$ dependent on the context
$(x_1, \ldots, x_T)$:
\[
\mathbf{h}_t = \phi(\mathbf{x}, t) \neq \phi(\mathbf{x}', t) \quad \text{if } \mathbf{x} \neq \mathbf{x}'
\]
even when $x_t = x'_t$.
\end{definition}

This contrasts with static embeddings where $\phi(w)$ is context-independent.

% ------------------------------------------------------------
\section{ELMo: Context via Bidirectional RNN}
\label{sec:pre:elmo}
% ------------------------------------------------------------

\begin{definition}[ELMo]
ELMo (Embeddings from Language Models, Peters et al., 2018) produces
contextual representations by combining layers of a BiLSTM pre-trained as a
bidirectional language model:
\[
\text{ELMo}_t = \gamma \sum_{\ell=0}^{L} s_\ell \mathbf{h}_t^\ell
\]
where $\mathbf{h}_t^\ell$ is the hidden state of layer $\ell$ at position $t$,
$s_\ell$ are task-specific learned weights, and $\gamma$ is a scaling factor.
\end{definition}

% ------------------------------------------------------------
\section{BERT: Bidirectional Encoder Representations from Transformers}
\label{sec:pre:bert}
% ------------------------------------------------------------

\subsection{Architecture}

BERT uses a Transformer encoder (bidirectional self-attention) with the
following configurations:

\begin{center}
\begin{tabular}{lcc}
\toprule
& \textbf{BERT-base} & \textbf{BERT-large} \\
\midrule
Layers $N$ & 12 & 24 \\
Dimension $d$ & 768 & 1024 \\
Heads $H$ & 12 & 16 \\
Parameters & 110M & 340M \\
\bottomrule
\end{tabular}
\end{center}

\subsection{Pre-training Objectives}

\begin{definition}[Masked Language Model (MLM)]
The MLM objective randomly masks $15\%$ of input tokens and trains the model
to predict them:
\[
\mathcal{L}_{\text{MLM}} = -\sum_{t \in \mathcal{M}} \log P(x_t \mid \mathbf{x}_{\setminus \mathcal{M}})
\]
where $\mathcal{M}$ is the set of masked positions. Among selected tokens:
$80\%$ are replaced with \texttt{[MASK]}, $10\%$ with a random token, $10\%$
unchanged.
\end{definition}

\begin{definition}[Next Sentence Prediction (NSP)]
The NSP objective trains the model to predict whether sentence $B$ follows
sentence $A$ in the original corpus:
\[
\mathcal{L}_{\text{NSP}} = -\log P(\text{IsNext} \mid \texttt{[CLS]}, A, \texttt{[SEP]}, B)
\]
\end{definition}

\begin{warning}[title=\textbf{MLM limitations}]
Masking creates a mismatch between pre-training (\texttt{[MASK]} tokens present)
and inference (no \texttt{[MASK]}). Subsequent work (RoBERTa, ALBERT) showed
that NSP is unnecessary and that training benefits from longer sequences and
more data.
\end{warning}

\subsection{WordPiece Tokenization}

BERT uses WordPiece, a BPE variant that maximizes likelihood:
\[
\text{score}(a, b) = \frac{\text{count}(ab)}{\text{count}(a) \cdot \text{count}(b)}
\]

% ------------------------------------------------------------
\section{GPT: Generative Pre-trained Transformer}
\label{sec:pre:gpt}
% ------------------------------------------------------------

\begin{definition}[GPT]
GPT (Radford et al., 2018) uses a Transformer decoder (causal self-attention)
pre-trained with an autoregressive objective:
\[
\mathcal{L}_{\text{GPT}} = -\sum_{t=1}^{T} \log P(x_t \mid x_1, \ldots, x_{t-1})
\]
\end{definition}

\begin{center}
\begin{tabular}{lccc}
\toprule
\textbf{Model} & \textbf{Parameters} & \textbf{Layers} & \textbf{Context} \\
\midrule
GPT-1 & 117M & 12 & 512 \\
GPT-2 & 1.5B & 48 & 1024 \\
GPT-3 & 175B & 96 & 2048 \\
GPT-4 & undisclosed & --- & 128K \\
\bottomrule
\end{tabular}
\end{center}

\begin{remark}
GPT-3 demonstrated emergent \emph{few-shot learning} capabilities: the ability
to solve new tasks from a few examples presented in the prompt, without weight
updates.
\end{remark}

% ------------------------------------------------------------
\section{T5: Text-to-Text Transfer Transformer}
\label{sec:pre:t5}
% ------------------------------------------------------------

\begin{definition}[T5]
T5 (Raffel et al., 2020) reformulates all NLP tasks as text-to-text problems.
The model is an encoder-decoder Transformer pre-trained with a denoising
(\emph{span corruption}) objective: contiguous token spans are replaced with
sentinels, and the decoder must reconstruct the masked spans.
\end{definition}

\begin{keyformulas}[title=\textbf{Architecture comparison}]
\begin{center}
\begin{tabular}{lccc}
\toprule
& \textbf{BERT} & \textbf{GPT} & \textbf{T5} \\
\midrule
Architecture & Encoder & Decoder & Enc-Dec \\
Attention & Bidirectional & Causal & Bid. + Causal \\
Objective & MLM + NSP & Autoregressive & Span corruption \\
Use case & Understanding & Generation & Both \\
\bottomrule
\end{tabular}
\end{center}
\end{keyformulas}

% ------------------------------------------------------------
\section{Scaling Laws}
\label{sec:pre:scaling}
% ------------------------------------------------------------

\begin{theorem}[Scaling Laws (Kaplan et al.)]
The test loss $L$ of a language model follows power laws with respect to the
number of parameters $N$, dataset size $D$, and compute budget $C$:
\begin{align}
L(N) &\propto N^{-\alpha_N} & \alpha_N &\approx 0.076 \\
L(D) &\propto D^{-\alpha_D} & \alpha_D &\approx 0.095 \\
L(C) &\propto C^{-\alpha_C} & \alpha_C &\approx 0.050
\end{align}
(for GPT-type models on English data).
\end{theorem}

\begin{remark}
The Chinchilla scaling laws (Hoffmann et al., 2022) show that optimal compute
budget allocation follows $N \propto D$: model size and data size should be
increased at the same rate.
\end{remark}

% ------------------------------------------------------------
\section{Using Hugging Face}
\label{sec:pre:hf}
% ------------------------------------------------------------

\begin{codeblock}[title=\textbf{BERT for text classification}]
\begin{minted}{python}
from transformers import AutoTokenizer, AutoModelForSequenceClassification
from transformers import pipeline

# Sentiment pipeline
clf = pipeline("sentiment-analysis",
               model="nlptown/bert-base-multilingual-uncased-sentiment")
result = clf("This NLP course is truly excellent!")
print(result)
# [{'label': '5 stars', 'score': 0.73}]

# Direct usage
tokenizer = AutoTokenizer.from_pretrained("bert-base-uncased")
model = AutoModelForSequenceClassification.from_pretrained(
    "bert-base-uncased", num_labels=2
)
inputs = tokenizer("This is a great course!", return_tensors="pt")
outputs = model(**inputs)
logits = outputs.logits  # (1, 2)
\end{minted}
\end{codeblock}

\begin{codeblock}[title=\textbf{GPT-2 for text generation}]
\begin{minted}{python}
from transformers import pipeline

generator = pipeline("text-generation", model="gpt2")
text = generator(
    "Natural language processing is",
    max_length=50,
    num_return_sequences=1,
    temperature=0.7
)
print(text[0]["generated_text"])
\end{minted}
\end{codeblock}

\begin{codeblock}[title=\textbf{T5 for summarization}]
\begin{minted}{python}
from transformers import pipeline

summarizer = pipeline("summarization", model="t5-base")
article = """
Natural language processing (NLP) is a subfield of linguistics,
computer science, and artificial intelligence concerned with the
interactions between computers and human language. The goal is to
enable computers to understand, interpret, and generate human language.
"""
summary = summarizer(article, max_length=30, min_length=10)
print(summary[0]["summary_text"])
\end{minted}
\end{codeblock}

% ------------------------------------------------------------
\section{Variants and Evolution}
\label{sec:pre:variants}
% ------------------------------------------------------------

\begin{description}
    \item[RoBERTa] (Liu et al., 2019): removes NSP, trains longer with more
          data and longer sequences.
    \item[ALBERT] (Lan et al., 2020): embedding factorization and cross-layer
          parameter sharing to reduce size.
    \item[DeBERTa] (He et al., 2021): disentangled attention with relative
          position.
    \item[LLaMA] (Touvron et al., 2023): Meta's open model family with
          RMSNorm, RoPE, and SwiGLU.
    \item[Mistral / Mixtral] (2023--2024): grouped-query attention (GQA),
          mixture of experts (MoE).
\end{description}

% ------------------------------------------------------------
\section{Feature Extraction and Probing}
\label{sec:pre:probing}
% ------------------------------------------------------------

\begin{definition}[Probing Classifiers]
\emph{Probing classifiers} are simple classifiers (logistic regression)
trained on frozen representations from a pre-trained model to evaluate what
linguistic information is captured by each layer.
\end{definition}

Studies have shown that lower layers capture morphology and syntax, while upper
layers encode more semantic information.

% ------------------------------------------------------------
\section{Exercises}
\label{sec:pre:exercises}
% ------------------------------------------------------------

\begin{exercise}[$\star$]
Compare the parameter counts of BERT-base and GPT-2 (117M). What architectural
differences explain the gaps?
\end{exercise}

\begin{exercise}[$\star$]
Explain why BERT cannot be directly used for autoregressive text generation.
\end{exercise}

\begin{exercise}[$\star\star$]
Implement a probing classifier to evaluate whether BERT layer-6 representations
encode POS tagging information. Use the Universal Dependencies dataset.
\end{exercise}

\begin{exercise}[$\star\star$]
Reproduce BERT's masking experiment: compute the pseudo-log-likelihood
perplexity
$\text{PLL}(\mathbf{x}) = \sum_{t=1}^{T} \log P(x_t \mid \mathbf{x}_{\setminus t})$
on a set of sentences.
\end{exercise}

\begin{exercise}[$\star\star\star$]
Study the impact of scaling laws: train Transformer language models of
increasing sizes (1M, 10M, 100M parameters) on the same corpus and plot the
loss vs.\ parameters curve on a log-log scale.
\end{exercise}
